\providecommand{\PacketTheoremNumber}{6.4}
\providecommand{\PacketBoundEquation}{6.8}
\providecommand{\PacketPhaseSection}{6.1}
\providecommand{\PacketFamilyDefinition}{6.3}
\providecommand{\PacketProofSection}{7}

\section{Introduction}\label{sec:introduction}

Let \(\Sigma\subset\R^3\) be a compact smooth surface, with induced
surface measure \(d\sigma\). Its Fourier extension operator is
\[
 E_\Sigma f(x)=\int_\Sigma f(\xi)e^{2\pi i x\cdot\xi}\dd\sigma(\xi).
\]
We allow smooth boundary. Positive curvature means that the second
fundamental form is definite at every point, with the normal chosen
on each chart so that the form is positive definite.

\begin{theorem}\label{thm:main}
Let \(\Sigma\subset\R^3\) be a compact smooth positively curved
surface, possibly with smooth boundary. For every \(3<p<\infty\),
there is a finite constant \(C_{\Sigma,p}\) such that
\begin{equation}\label{eq:main}
 \|E_\Sigma f\|_{L^p(\R^3)}
 \le C_{\Sigma,p}\|f\|_{L^p(\Sigma)}
\end{equation}
for every complex-valued \(f\in L^p(\Sigma)\).
\end{theorem}

The theorem establishes the diagonal extension conjecture in three
dimensions for positively curved surfaces. It includes the sphere,
compact elliptic paraboloid patches, and general smooth elliptic
graphs. The surface and the exponent are fixed throughout each
estimate.

\paragraph{Background.}
The restriction problem originates in Stein's work; the
Tomas--Stein theory gives \(L^2(\Sigma)\to L^4(\R^3)\)
for positively curved surfaces \cite{Tomas1975,Stein1979}.
The bilinear framework of Tao, Vargas, and Vega
\cite{TaoVargasVega1998}, Wolff's cone estimate \cite{Wolff2001},
and Tao's paraboloid estimate \cite{Tao2003} developed the role of
wave packets and transverse interactions. Bennett, Carbery, and Tao's
multilinear estimates \cite{BennettCarberyTao2006} and
Bourgain--Guth's approach \cite{BourgainGuth2011} further connected
these interactions to the linear problem.

Guth's polynomial-partitioning argument
established the bounded-data estimate for \(p>13/4\)
\cite[Theorem~0.1]{Guth2016}. Shayya obtained diagonal estimates in that range
\cite[Corollary~2.1(iii)]{Shayya2017}; see also Kim's treatment
\cite[Theorem~1.2 and the following remark]{Kim2017}. Wang's broom
method further improved bounded-data restriction for the truncated
paraboloid \cite[Theorem~1.4]{Wang2022Brooms}. Wang and Wu proved
the diagonal bound for \(p>22/7\) on compact positively curved
surfaces, including surfaces with boundary, using decoupling and
two-ends Furstenberg inequalities \cite[Theorem~0.2]{WangWu2024}.

For the sphere and quadratic graphs, factorization gives a separate
route from bounded-data estimates to diagonal estimates with an
arbitrarily small increase in the exponent; see
Bourgain \cite{Bourgain1991Besicovitch} and Buschenhenke
\cite[Theorem~1.1 and Remark~1.6]{Buschenhenke2024}. The argument here treats general smooth
elliptic charts through the packet theorem's full array formulation
and the explicit phase continuation. The packet and maximal inputs
play distinct roles: the former controls oscillatory coefficients,
while the latter controls their positive tube density.

The compact paraboloid case also gives a consequence for free
Schr\"odinger evolution. Let \(W^{s,p}(\R^2)\) denote the inhomogeneous
Bessel-potential space with norm
\(\|f\|_{W^{s,p}}=\|(I-\Delta)^{s/2}f\|_{L^p(\R^2)}\).

\begin{corollary}[Free Schr\"odinger local smoothing]\label{cor:schrodinger-smoothing}
For every \(3<p<\infty\) and \(s>2-6/p\), the free Schr\"odinger
propagator, initially defined on Schwartz functions, extends boundedly
from \(W^{s,p}(\R^2)\) to \(L^p(\R^2\times[0,1])\):
\begin{equation}\label{eq:schrodinger-smoothing}
 \|e^{it\Delta}f\|_{L^p(\R^2\times[0,1])}
 \le C_{p,s}\|f\|_{W^{s,p}(\R^2)}.
\end{equation}
\end{corollary}

The spacetime norm is global in space and local in time. Rogers's necessary
condition \cite[Section~2]{Rogers2008} shows that the Sobolev threshold
cannot be lowered. We deduce the corollary from the compact-frequency
transfer of Lee, Rogers, and Seeger \cite[Theorem~1.1]{LeeRogersSeeger2013}
in Section~\ref{sec:surfaces}.

A second consequence concerns local oscillatory integrals with phase
\(N[x\cdot y+\psi(t;y)]\), where \(x,y\in\R^2\), \(t\in\R\),
and the amplitude is smooth and compactly supported. Suppose
\(D_y^2\partial_t\psi\) is definite and its time derivative is a
scalar multiple of itself. A change of output variables then reduces
the phase to a fixed elliptic graph. Corollary~\ref{cor:bourgain-oscillatory}
gives the resulting \(L^p\) estimate with the scaling factor
\(N^{-3/p}\) for every finite \(p>3\). The additional phase condition
is the translation-invariant form of Bourgain's condition; we explain
the straightening directly in Section~\ref{sec:surfaces}, following
the setting of Gao--Liu--Xi \cite{GaoLiuXi2025}.

The proof of Theorem~\ref{thm:main} uses two theorem inputs. The packet
propagation theorem \cite[Theorem~\PacketTheoremNumber]{OpenAIPacket2026}
controls
arbitrary finite complex coefficient arrays through an elliptic
capacity, with coordinate masks permitted. The Kakeya maximal
theorem \cite[Theorem~1.1]{OpenAIKakeya2026} controls nonnegative
weighted tube concentration. We state both inputs in the forms used
here and verify their hypotheses. The remainder of the proof gives
the density conversion, the estimates for rough \(L^3\) graph data,
and the passage to general surfaces.

\paragraph{The argument.}
After fixing a graph chart, an explicit continuation places its
phase in the class allowed by the packet theorem. The construction
matches the quadratic Taylor data in a transition annulus of a
continued spherical patch and preserves uniform ellipticity.
It applies equally to boundary charts, where the graph extends
smoothly and the datum is extended by zero.

The packet theorem gives a cubic sample estimate with a small
spatial power loss and an elliptic-capacity factor. A decomposition
according to the capacity of the remaining coefficients allows this
factor to be summed using the positive weighted tube estimate.
The latter follows from the Kakeya maximal theorem by duality;
its direction and position formulation follows the classical
weighted tube framework of Tao, Vargas, and Vega
\cite[arXiv version~1, Section~3]{TaoVargasVega1998}.

To pass to \(L^3\) data, we prove one joint estimate for a sparse
family of balls. Removing a small union of frequency resonances
makes the corresponding packet analyses almost orthogonal. The
bound is uniform in the ball locations and in arbitrarily large
separations. A covering of each level set then removes the spatial
loss. This follows Tao's sparse-ball strategy for epsilon removal
\cite{Tao1999Epsilon};
see the detailed treatments in Bourgain--Guth's appendix
\cite{BourgainGuth2011} and Kim's appendix \cite{Kim2017}.
We prove the precise covering and summation needed for the present
estimate.

The resulting graph bound has \(L^3\) input and \(L^p\) output
for every \(p>3\). A finite chart decomposition, with the surface
and physical-space Jacobians retained, gives
Equation~\eqref{eq:main} by the finite-measure embedding
\(L^p(\Sigma)\subset L^3(\Sigma)\).

Figure~\ref{fig:argument-map} records the two inputs and the successive
reductions.

\begin{figure}[tbp]
\centering
\begin{tikzpicture}[
  font=\small,
  amap box/.style={
    draw=black!60, fill=black!2, rounded corners=2pt,
    align=center, text width=.37\textwidth,
    inner xsep=7pt, inner ysep=6pt, minimum height=12mm
  },
  amap input/.style={amap box, draw=blue!45!black, fill=blue!5},
  amap result/.style={amap box, text width=.53\textwidth},
  amap arrow/.style={-{Stealth[length=2mm]}, draw=black!65, semithick}
]
\node[amap box] (phase) at (-.235\textwidth,0)
  {Positively curved chart\\Insertion into an allowed phase};
\node[amap input] (kakeya) at (.235\textwidth,0)
  {$L^3$ Kakeya maximal theorem\\Loss $\rho^{-s}$ for each $s>0$};

\node[amap input, below=6mm of phase] (packet)
  {Packet theorem\\Complex arrays; coordinate masks};
\node[amap box, below=6mm of kakeya] (tubes)
  {Weighted tube estimate\\Nonnegative tube weights};

\coordinate (join) at ($(packet.south)!.5!(tubes.south)$);
\node[amap result, below=9mm of join] (samples)
  {Density grouping\\Sampled cubic estimate};
\node[amap result, below=6mm of samples] (sparse)
  {Joint estimate on sparse balls\\Compactly supported $L^3$ data};
\node[amap result, below=6mm of sparse] (global)
  {Level-set cover and loss removal\\Graph $L^3\to L^p$, $p>3$};
\node[amap result, below=6mm of global] (surface)
  {Finite charts and surface measure\\Surface $L^p\to L^p$, $p>3$};

\draw[amap arrow] (phase) -- (packet);
\draw[amap arrow] (kakeya) -- (tubes);
\draw[amap arrow] (packet.south) -- ([xshift=-.15\textwidth]samples.north);
\draw[amap arrow] (tubes.south) -- ([xshift=.15\textwidth]samples.north);
\draw[amap arrow] (samples) -- (sparse);
\draw[amap arrow] (sparse) -- (global);
\draw[amap arrow] (global) -- (surface);
\end{tikzpicture}
\caption{The two independent inputs meet in the sampled cubic estimate.
The sparse-ball bound then supplies the common $L^3$ control needed for
globalization and passage to the surface.}
\label{fig:argument-map}
\end{figure}


\paragraph{Dependence on the exponent.}
Let \(C_{\Sigma,p}^{\mathrm{best}}\) denote the operator norm in
Equation~\eqref{eq:main}. Write \(M(s)\ge1\) for the constant in
the Kakeya maximal estimate with loss \(\delta^{-s}\), and let
\(P_\Sigma(\kappa,\eta)\ge1\) be the largest packet constant
among the finitely many fixed chart and window setups used below.
For \(3<p\le4\), put \(\theta=10^{-6}(p-3)^2\). The proof gives
\begin{equation}\label{eq:quantitative-intro}
 C_{\Sigma,p}^{\mathrm{best}}
 \le C(\Sigma)(p-3)^{-7/p}
 P_\Sigma(\theta,\theta)^{1/p}M(\theta)^{3/(2p)},
\end{equation}
where \(C(\Sigma)\) is independent of \(p\in(3,4]\).
The two input theorems supply finite constants at each positive
loss exponent. No rate for their growth is assumed, so
Equation~\eqref{eq:quantitative-intro} records the elementary
losses without asserting a polynomial bound in \((p-3)^{-1}\).
For nonempty \(\Sigma\), a fixed smooth datum also gives
\begin{equation}\label{eq:lower-intro}
 C_{\Sigma,p}^{\mathrm{best}}
 \ge c(\Sigma)(p-3)^{-1/p},\qquad 3<p\le4.
\end{equation}
Proposition~\ref{prop:lower} proves this comparison.

\paragraph{Conventions and organization.}
Write \(e(t)=\exp(2\pi i t)\). All functions and coefficient
arrays may be complex-valued. Constants may depend on the fixed
surface, charts, frequency supports, and windows; dependence on
the small loss exponents is displayed through \(P\) and \(M\).
Section~\ref{sec:setup} fixes the local graph problem. The next
sections give the packet interface, the density conversion, and
the joint sparse-ball estimate. The global argument tracks the
dependence on \(p\), and Section~\ref{sec:surfaces} completes
the transfer to \(\Sigma\) and proves the local-smoothing and
oscillatory-integral consequences.
